\section{Formalisation, provenance and tool disclosure}\label{sec:disclosure}

\paragraph{Machine-checked part.} The finite algebra behind \Cref{sec:windows} and the core of \Cref{thm:nominorant} are
formalised in Lean~4 (toolchain \texttt{v4.26.0}) against Mathlib \cite{Mathlib} (tag \texttt{v4.26.0}, revision \texttt{2df2f0150c27}).
\Cref{app:lean} lists the thirteen files, their declaration counts, and the statements they cover; the last of them formalises the
core of the obstruction theorem (\Cref{lem:indep} and the Fatou--independence argument of \Cref{thm:nominorant}) with the cutoff
budget of \Cref{lem:budget} as an explicit hypothesis. Every file compiles with
\texttt{\#print axioms} reporting only \texttt{propext}, \texttt{Classical.choice}, \texttt{Quot.sound}; no project axiom is used.
What is formalised is exactly what is finite: curvature identities and the explicit product bound for finite transforms, lattice
identities, second-mode overlap, the sign-free Ritz bookkeeping, the decomposition $K=\Gamma-c_LI-2\beta\beta^*$ of the window
matrix, the von Mangoldt divisibility energy, and the finite obstructions (tradeoff, eigen-tradeoff, Schur, signed-PSD plant).
The analytic statements (\Cref{thm:FPhi,prop:radical,thm:GS,lem:budget}) are paper theorems; the explicit formula \eqref{eq:EF}
is imported from the literature and is not a project axiom.

\paragraph{Numerical part.} All finite-window quantities were computed in interval arithmetic (\texttt{arb}, 220--300 decimal digits) with the mode truncation chosen from the required relative precision $\approx\sqrt{\lambda_1}$; float
evaluations were used only for figures. The identities \eqref{eq:EF} on $f_0$ and on plane waves, the masses \eqref{eq:masses} and the
constants in \Cref{app:constants} were computed with \texttt{mpmath} at 30--50 digits, with the first $3000$ zeros of $\zeta$ where
zeros enter.

\paragraph{Use of automated reasoning tools.} Following the Leiden Declaration on AI and Mathematics as summarised in
\cite[\S8]{TERENCETAO-2026}, we disclose that large language models were used throughout: as adversarial reviewers of proof
sketches, as drafters of Lean code that was then checked by the kernel, and as literature scouts whose every claim was verified against
the primary source or discarded. Such a reviewer first wrote several of the proofs in \Cref{sec:obstruction}, in particular \Cref{lem:budget,thm:nominorant},
in response to a formally specified request. A second, blind agent then re-derived them independently from the definitions,
and the author checked them numerically. Predictions about the outcome of each such round were
registered before the round and are kept in the project log; the author's own predictions were wrong more than once (for instance,
a probability of $0.40$ was assigned to \Cref{thm:nominorant} before it was proved). Credit and responsibility for every statement here
remain with the author.

\paragraph{Relation to earlier work of the author.} The preprint \cite{Malamutmann2025} attempted to prove positivity of the same form by
separate archimedean and prime estimates on a wide cone of tests. Item~1 of \Cref{sec:ledger} shows why that cannot work; the present paper supersedes it.
